\section*{PHYSIQUE (13 points)}
\section{Ondes ultrasonores}
Les ondes ultrasonores sont des ondes mécaniques qui peuvent se propager dans
des milieux différents. Elles engendrent dans des conditions bien définies
certains phénomènes physiques.

\vskip2mm
\begin{minipage}{0.59\linewidth}
    Pour déterminer la célérité d'une onde ultrasonore de fréquence $N$ dans
    deux milieux différents, on utilise un dispositif constitué d'un émetteur E
    et d'un récepteur R fixés aux extrémités d'un tube.
    E et R sont reliés à un oscilloscope.

    \begin{donnees}
        \begin{itemize}[itemsep=-5pt]
            \item Distance émetteur-récepteur~: $D=\mathrm{ER}=\qty{1}{\m}$~;
            \item $N=\qty{40}{\kHz}$.
        \end{itemize}
    \end{donnees}
\end{minipage}
\hfill
\begin{minipage}{0.40\linewidth}
    \flushright
    \includegraphics[width=\textwidth]{2025_05_09_e1619f2897d39ae202a1g-4(1)}
\end{minipage}


\begin{enumerate}
    \item L'onde ultrasonore est-elle une onde longitudinale ou transversale~?
    \item On remplit le tube par de l'eau. L'oscillogramme ci-contre représente
          le signal émis par E et celui reçu par R.

          Recopier sur votre copie le numéro de la question et écrire la lettre
          correspondante à la proposition vraie.
          \begin{enumerate}
              \item La célérité des ultrasons dans l'eau vaut~:
                    \begin{table}[H]
                        \flushright
                        \setlength{\arrayrulewidth}{1pt}
                        \renewcommand{\arraystretch}{1.5}
                        \begin{tabularx}{.90\textwidth}{|p{3mm}|X|p{3mm}|X|p{3mm}|X|p{3mm}|X|}
                            \hline
                            A & $c=\qty{1520}{\m\per\s}$ &
                            B & $c=\qty{620}{\m\per\s}$ &
                            C & $c=\qty{1667}{\m\per\s}$ &
                            D & $c=\qty{330}{\m\per\s}$ \\
                            \hline
                        \end{tabularx}
                    \end{table}
              \item La longueur d'onde de l'onde ultrasonore vaut~:
                    \begin{table}[H]
                        \flushright
                        \setlength{\arrayrulewidth}{1pt}
                        \renewcommand{\arraystretch}{1.5}
                        \begin{tabularx}{.90\textwidth}{|p{3mm}|X|p{3mm}|X|p{3mm}|X|p{3mm}|X|}
                            \hline
                            A & $\lambda=\qty{25.2}{\mm}$ &
                            B & $\lambda=\qty{30.5}{\mm}$ &
                            C & $\lambda=\qty{37.2}{\mm}$ &
                            D & $\lambda=\qty{41.7}{\mm}$ \\
                            \hline
                        \end{tabularx}
                    \end{table}
          \end{enumerate}

    \item On remplace l'eau par un autre liquide, on constate que le décalage
          horaire entre le signal émis et le signal reçu est
          $\Delta t=\qty{0.9}{\s}$.

          La célérité des ultrasons dans le liquide, a-t-elle augmenté ou
          diminué par rapport à celle dans l'eau~? Justifier.
\end{enumerate}



\section*{Exercice 2 (5 points) : Evolution d'un système électrique}
Le comportement d'un système électrique dépend des éléments qui le constituent (Condensateur, Bobine,...). Selon les conditions initiales, l'évolution d'un tel système peut être décrite à l'aide de certains paramètres et grandeurs électriques ou énergétiques.

\section*{Partie 1 : Détermination de la capacité d'un condensateur}
On réalise la charge d'un condensateur de capacité $C$, à l'aide d'un générateur idéal de courant qui débite un courant d'intensité constante $I_{0}=0,5 \mu A($ figure 1 ).

\begin{center}
\begin{tabular}{|l|l|}
\hline
Figure (1) & \begin{tabular}{l}
\includegraphics[width=\textwidth]{2025_05_09_e1619f2897d39ae202a1g-4}
 \\
Figure (2) \\
\end{tabular} \\
\hline
\end{tabular}
\end{center}

À l'instant $t_{0}=0$, on ferme l'interrupteur $K$. La figure (2) de la page $4 / 6$ représente l'évolution de la tension $u_{C}(t)$ aux bornes du condensateur.

    \begin{enumerate}
      \item Recopier sur votre copie le numéro de la question et écrire la lettre correspondante à la proposition vraie.\\
L'expression de $u_{C}$ est :\\
A $\quad u_{C}=\frac{C}{I_{0}} . t \left\lvert\, \mathbf{B} \quad u_{C}=\frac{I_{0}}{C} . t \quad \mathbf{C} \quad u_{C}=I_{0} . C . t \quad \mathbf{D} \quad u_{C}=C . t\right.$
      \item Vérifier que $C=0,5 \mu F$.
    \end{enumerate}
Partie 2: Étude de la décharge d'un condensateur à travers une bobine\\
À l'instant $t_{0}=0$, on branche le condensateur précédemment chargé aux bornes d'une bobine d'inductance $L$ et de résistance négligeable.

    \begin{enumerate}
      \item Établir l'équation différentielle vérifiée par la charge $q(t)$ du condensateur.
      \item La courbe de la figure (3), représente l'évolution de $q(t)$.\\
2.1. Nommer le régime d'oscillations que montre le graphe de la figure (3).\\
2.2. La solution de l'équation différentielle s'écrit : $q(t)=Q_{m} \cdot \cos \left(\frac{2 \pi}{T_{0}} \cdot t+\varphi\right)$.\\
\includegraphics[width=\textwidth]{2025_05_09_e1619f2897d39ae202a1g-5(1)}
    \end{enumerate}
Figure (3)\\
2.2.1. En exploitant le graphe de la figure (3), déterminer les valeurs de $Q_{m}, T_{0}$ et $\varphi$.\\
2.2.2. Calculer la valeur de $L$.\\
2.3. Expliquer qualitativement la conservation de l'énergie totale du circuit ( $L C$ ) et calculer sa valeur.\\
2.4. Déterminer la valeur maximale de l'intensité du courant dans le circuit.

\section{Évolution d'un système mécanique}
Les mouvements des systèmes mécaniques dépendent de la nature des actions
mécaniques qui leurs sont appliquées.

L'étude de l'évolution temporelle de ces systèmes permet de déterminer certaines
grandeurs dynamiques et cinématiques et d'expliquer certains aspects
énergétiques.

Cet exercice vise l'étude du mouvement de translation rectiligne d'un solide sur
un plan incliné et l'étude du mouvement du système oscillant \{solide -
ressort\}.

Dans cet exercice tous les frottements sont supposés négligeables.

\subsection{Mouvement d'un solide sur un plan incliné}
\begin{minipage}{0.45\linewidth}
    On considère un solide ($S$) de masse $m$ susceptible de glisser selon la
    ligne de plus grande pente d'un plan incliné faisant un angle $\alpha$ avec
    l'horizontal.
    
    Le solide $(S)$ démarre sans vitesse initiale, à l'instant $t_{0}=0$ à
    partir de la position $O$ sous l'action d'une force motrice $\vec{F}$
    constante.
    
    Le solide ($S$) passe par la position $A$ avec la vitesse $v_{A}$. On étudie
    le mouvement du centre d'inertie $G$ du solide $(S)$ dans un repère $(O,
    \vec{i})$ lié à la Terre supposé galiléen (figure 1 ).
    
    L'abscisse de $G$ à $t_{0}=0$ est $x_{G}=x_{0}=0$.
\end{minipage}
\hfill
\begin{minipage}{0.45\linewidth}
    \begin{figure}[H]
        \centering
        \includegraphics[width=\textwidth]{2025_05_09_e1619f2897d39ae202a1g-5}
        \caption{}
        \label{fig:}
    \end{figure}
\end{minipage}
\begin{donnees}
    \begin{itemize}
        \item $m=\qty{100}{\gram}$~; $g=\qty{10}{\a}$~; $\alpha=\ang{30}$~;
              $v_{A}=\qty{2,4}{\v}$.
    \end{itemize}
\end{donnees}

\begin{figure}[H]
    \centering
    \includegraphics[width=\textwidth]{2025_05_09_e1619f2897d39ae202a1g-6(1)}
    \caption{}
    \label{fig:}
\end{figure}


\begin{enumerate}
    \item En appliquant la deuxième loi de Newton, montrer que l'équation
          différentielle vérifiée par $x_{G}$ s'écrit~:
          $\frac{\mathrm{d}^{2}x_{G}}{\mathrm{d}t^{2}}=\frac{F}{m}-g\sin\alpha$.
    \item La figure~(2) donne l'évolution de la vitesse $v(t)$.
          \begin{enumerate}
              \item Déterminer graphiquement la valeur de l'accélération du
                    mouvement de $G$.
              \item Calculer l'intensité de la force $\vec{F}$.
          \end{enumerate}
    \item À partir de la position $A$, le solide ($S$) n'est plus soumis à la
          force motrice $\vec{F}$ et s'arrête en une position $B$.

          On choisit $A$ comme nouvelle origine des abscisses et l'instant de
          passage de $G$ par A comme nouvelle origine des dates.
          \begin{enumerate}
              \item En utilisant l'équation différentielle établie dans la
                    question~(1), montrer que le mouvement de $G$ entre $A$ et
                    $B$ est rectiligne uniformément varié.
              \item Déterminer la distance $AB$.
          \end{enumerate}
\end{enumerate}

\subsection{Mouvement d'un système \{solide - ressort\}}
\begin{minipage}{0.65\linewidth}
    On considère le système \{solide ($S$)-ressort\} représenté sur la
    figure~(\ref{fig:solide_ressort}). Le ressort est à spires non jointives,
    d'axe horizontal, de masse négligeable et de raideur $K$. On étudie le
    mouvement du centre d'inertie $G$ du solide ($S$) de masse
    $m=\qty{100}{\gram}$ dans un repère $(O,\ex)$ lié à la Terre supposé
    galiléen.
    
    À l'équilibre $x_{G}=x_{0}=0$.
\end{minipage}
\hfill
\begin{minipage}{0.30\linewidth}
    \begin{figure}[H]
        \centering
        \begin{tikzpicture}
    \node[inner sep=0pt,minimum width=4.5cm,minimum height=3mm,
        top color=black!70] (sol) {};
    % Axe xx'
    \coordinate (O) at ($(sol.south west)!.8!(sol.south east)$);
    \draw[-Latex] (sol.south west)
        node[below right,inner xsep=0pt] {$x^{\prime}$} -- (O)
        node[below=0.4mm] {$O$} -- (sol.south east) -- ++(0.7,0)
        node[below left,inner xsep=0pt,inner ysep=2mm] {$x$};
    % Fixed
    \node[inner sep=0pt,draw,minimum width=3mm,minimum height=7mm,
        anchor=south,pattern=north west lines,
        thick] (fix) at ([xshift=3mm,yshift=0.4pt]sol.north west) {};
    % Solide (S)
    \node[inner sep=0pt,draw,minimum width=7mm,minimum height=7mm,
        anchor=south,rounded corners=2pt,thick,label=$\mathrm{(S)}$,
        fill=lightgray] (S) at
        ([yshift=0.4pt]O |- sol.north) {};
        %([xshift=-6mm,yshift=0.4pt]sol.north east) {};
    \node[inner sep=1pt,circle,fill,
        label={[right=-0.8mm,font=\scriptsize]:$G$}] (G) at (S.center) {};
    \draw[{Rectangle[length=1pt,width=4pt]}-{Stealth[length=5pt]},
        thick] (O) -- ++(0.8,0)
        node[below left,inner xsep=0pt] {$\ex$};
    % Ressort
    \draw[decoration={
          coil,
          aspect=0.3,
          segment length=2.0mm,
          amplitude=2.4mm,
          pre length=1mm,
          post length=1mm},
          decorate,thick,black] (fix.east) -- (S.west);
\end{tikzpicture}

        \caption{}
        \label{fig:solide_ressort}
    \end{figure}
\end{minipage}

On écarte $(S)$ de sa position d'équilibre d'une distance $X_{m}$ et on
l'abandonne sans vitesse initiale à l'instant $t_{0}=0$. Le solide $(S)$
effectue 10 oscillations pendant la durée $\Delta t=\qty{3,14}{\second}$.
\begin{enumerate}
    \item Déterminer la valeur de la période propre $T_{0}$.
    \item Déduire la valeur de $K$.
    \item 
          \begin{minipage}[t][][t]{0.65\linewidth}
              On choisit l'état où le ressort n'est pas déformé comme référence
              de l'énergie potentielle élastique $E_{\mathrm{pe}}$ et le plan
              horizontal contenant $G$ comme état de référence de l'énergie
              potentielle de pesanteur $E_{\mathrm{pp}}$.

              La courbe de la figure (\ref{fig:Epe_de_x}) représente le
              diagramme d'énergie potentielle élastique $E_{\mathrm{pe}}=f(x)$.
              
              En exploitant le diagramme, déterminer les valeurs de~:
              \begin{enumerate}
                  \item L'amplitude $X_{m}$.
                  \item L'énergie mécanique $E_{m}$ du système oscillant.
                  \item La vitesse maximale du mouvement de ($S$).
              \end{enumerate}
          \end{minipage}
          \hfill
          \begin{minipage}[t][][t]{0.33\linewidth}
              \begin{figure}[H]
                  \centering
                  \begin{tikzpicture}
                      \def\K{40}
                      \begin{axis}[
                              xmin=-0.05,xmax=0.05,
                              ymin=0,ymax=38,
                              width=7cm,height=7cm,
                              xlabel=$x~(\unit{\meter})$,
                              ylabel=$E_{pe}~(\unit{\mJ})$,
                              every axis y label/.append style={
                              	at={(0.55,0.95)},anchor=west,
                              },
                              scaled ticks=false,
                              xtick distance=0.02, ytick distance=8,
                              xticklabels={,,,,{0,02}}, yticklabels={,,{8}},
                          ]
                          \addplot[domain=-0.04:0.04,samples=200,
                              very thick,blue] {0.5*\K*x^2*1e3}
                              coordinate[pos=0](A) coordinate[pos=1](B);
                      \end{axis}
                      \foreach \n/\a in {A/160,B/20}{%
                      \draw[very thick,blue]
                          ([yshift=-5mm]\n) -- ([yshift=5mm]\n);
                      \foreach \f in {0.02,0.14,...,1}
                      \draw[very thick,blue]
                        ($([yshift=-5mm]\n)!\f!([yshift=5mm]\n)$) -- ++(\a:.2);
                      }
                  \end{tikzpicture}
                  \caption{}
                  \label{fig:Epe_de_x}
              \end{figure}
          \end{minipage}
\end{enumerate}


